% 第1章 1.5.1：两种测度下的包含方向（示意图，嵌套框表示集合包含）。
% (a) (0,1] 上 Lebesgue 测度（有限总质量）：L^∞ ⊂ L^2 ⊂ L^1；x^{-1/4}∈L^2\L^∞，x^{-1/2}∈L^1\L^2。
% (b) N 上计数测度：ℓ^1 ⊂ ℓ^2 ⊂ ℓ^∞；(1/k^2)∈ℓ^1，(1/k)∈ℓ^2\ℓ^1，(1,1,...)∈ℓ^∞\ℓ^2。
\begin{tikzpicture}[font=\figurefont\footnotesize,
  box/.style={draw=#1,line width=.75pt,rounded corners=1.8pt},
  lab/.style={font=\figurefont\footnotesize,anchor=north west,inner sep=1.2pt},
  ex/.style={font=\figurefont\scriptsize,text=FigureInk},dot/.style={circle,fill=FigureStroke,inner sep=0pt,minimum size=3pt}]
\begin{scope}
\node[font=\figurefont\footnotesize,anchor=south] at (2.4,3.15) {(a) $(0,1]$，Lebesgue测度};
\draw[box=FigureYellow,fill=FigureYellowFill] (0,0) rectangle (4.8,3.0);
\draw[box=FigureGreen,fill=FigureGreenFill] (0.25,0.2) rectangle (4.55,2.15);
\draw[box=FigureBlue,fill=FigureBlueFill] (0.5,0.4) rectangle (4.3,1.3);
\node[lab,text=FigureYellow] at (0,3.0) {$L^1$};
\node[lab,text=FigureGreen] at (0.25,2.15) {$L^2$};
\node[lab,text=FigureBlue] at (0.5,1.3) {$L^\infty$};
\node[dot] at (2.0,2.55) {}; \node[ex,anchor=west] at (2.1,2.55) {$x^{-1/2}$};
\node[dot] at (2.0,1.72) {}; \node[ex,anchor=west] at (2.1,1.72) {$x^{-1/4}$};
\node[dot] at (2.0,0.82) {}; \node[ex,anchor=west] at (2.1,0.82) {$1$, $x$};
\end{scope}
\begin{scope}[xshift=5.6cm]
\node[font=\figurefont\footnotesize,anchor=south] at (2.4,3.15) {(b) $\mathbb N$，计数测度};
\draw[box=FigureYellow,fill=FigureYellowFill] (0,0) rectangle (4.8,3.0);
\draw[box=FigureGreen,fill=FigureGreenFill] (0.25,0.2) rectangle (4.55,2.15);
\draw[box=FigureBlue,fill=FigureBlueFill] (0.5,0.4) rectangle (4.3,1.3);
\node[lab,text=FigureYellow] at (0,3.0) {$\ell^\infty$};
\node[lab,text=FigureGreen] at (0.25,2.15) {$\ell^2$};
\node[lab,text=FigureBlue] at (0.5,1.3) {$\ell^1$};
\node[dot] at (1.7,2.55) {}; \node[ex,anchor=west] at (1.8,2.55) {$(1,1,1,\ldots)$};
\node[dot] at (1.7,1.72) {}; \node[ex,anchor=west] at (1.8,1.72) {$(1/k)_{k\ge1}$};
\node[dot] at (1.7,0.82) {}; \node[ex,anchor=west] at (1.8,0.82) {$(1/k^2)_{k\ge1}$};
\end{scope}
\end{tikzpicture}
